\section{Installation der elektrischen Komponenten}
Die elektrische Installation basiert auf der Versorgungsspannung von +5~V.  
Alle LED-Streifen sowie der Mikrocontroller werden damit versorgt und haben einen gemeinsamen Ground.  

Die \gls{led}-Ansteuerung ist in drei separate \gls{led}-Stränge unterteilt, wobei jeder Strang aus mehreren unterschiedlich langen \gls{led}-Streifen besteht, die durch Lötverbindungen miteinander verbunden sind, um Verbindungen über Kurven zu realisieren.
Der erste Strang hat 256 \gls{ws2812b}-\gls{led}s und wird für die Anzeige von Stunden und Minuten, der Außentemperatur und dem Wetter verwendet.  
Der zweite Strang hat 60 \gls{ws2812b}-\gls{led}s und wird für die Sekundenanzeige genutzt.
Der dritte Strang hat 38 \gls{ws2812b}-\gls{led}s und wird für die Anzeige der Zusatzfunktionen genutzt.
Die Datenleitungen der \gls{led}-Stränge werden jeweils auf separate \gls{gpio}-Pins des Mikrocontrollers geführt.  

Die Dimensionierung des Netzteils erfolgt auf Basis der maximal möglichen Leistungsaufnahme.  
Die Gesamtleistung ergibt sich aus der maximalen Anzahl an LEDs, die gleichzeitig aktiv sein können und deren spezifischer Leistungsaufnahme.  
Für die Berechnung kann die Beziehung  
\[
P_{\text{gesamt}} = N_{\text{LED}} \cdot P_{\text{LED}}
\]

verwendet werden.
Die Leistungsaufnahme pro LED kann über die Streckenleistung bestimmt werden welche im mitgelieferten Datenblatt der verwendeten \gls{led}-Streifen angegeben ist~\cite{AZDelivery}.  
\[
P_{\text{LED}} = \frac{P_{\text{pro Meter}}}{N_{\text{LED pro Meter}}}
\]

Für den verwendeten LED-Streifen ergibt sich mit der Angabe von  
\[
P_{\text{pro Meter}} = 6\,\text{W/m}
\]
und einer Bestückung mit  
\[
N_{\text{LED pro Meter}} = 30\,\text{LED/m}
\]

die Leistungsaufnahme einer einzelnen LED zu  
\[
P_{\text{LED}} = \frac{6\,\text{W/m}}{30\,\text{LED/m}}
= 0{,}2\,\text{W}.
\]

Im ungünstigsten Fall sind während des Boot-Modus kurzzeitig alle 256 LEDs des Hauptstreifens, bei voller Helligkeit, gleichzeitig aktiv. Daraus ergibt sich die maximale Gesamtleistung zu  
\[
P_{\text{gesamt}} = 256 \cdot 0{,}2\,\text{W}
= 51{,}2\,\text{W}.
\]

Die Stromaufnahme 
\[
I_{\text{gesamt}} = \frac{P_{\text{gesamt}}}{U}
\]

berechnet sich mit einer Versorgungsspannung von  
\[
U = 5\,\text{V}
\]
zu  
\[
I_{\text{gesamt}} = \frac{51{,}2\,\text{W}}{5\,\text{V}}
= 10{,}24\,\text{A}.
\]

Die mindestens erforderliche Netzteilleistung ergibt sich anschließend zu  
\[
P_{\text{Netzteil}} = U \cdot I_{\text{gesamt}}.
\]

Somit ergibt sich eine erforderliche Mindestleistung von 
\[
P_{\text{Netzteil}} = 5\,\text{V} \cdot 10{,}24\,\text{A}
= 51{,}2\,\text{W}.
\]

Obwohl im Boot-Modus kurzzeitig alle 256 LEDs des Hauptstreifens bei voller Helligkeit aktiviert werden, tritt dieser Zustand nur für etwa eine Sekunde auf. 
Für die thermische und dauerhafte Auslegung des Netzteils ist dieser kurzzeitige Spitzenwert nicht maßgebend.
Im regulären Betrieb liegt die Anzahl gleichzeitig aktivierter \gls{led}s bei 104 und damit deutlich unter diesem theoretischen Maximalwert.
Ein kurzzeitiges und geringfügiges Überschreiten der errechneten Leistungsaufnahme während des Boot-Modus kann vom Netzteil als Lastspitze bereitgestellt werden.

Aus diesem Grund wird ein Netzteil mit einer Ausgangsstromstärke von \SI{10}{\ampere} bei \SI{5}{\volt} verwendet.
Mit diesem ist auch die Leistungsaufnahme des \gls{esp32c6} und des \gls{dht22} Sensors mit genügend Reserve abgedeckt.